\documentclass[11pt,a4paper]{article}

\usepackage[utf8]{inputenc}
\usepackage[T1]{fontenc}
\usepackage{geometry}
\geometry{margin=1in}

% Paragraph spacing
\setlength{\parskip}{0.6\baselineskip plus 2pt minus 1pt}
\setlength{\parindent}{0pt}

\usepackage{mathtools,amsmath,amssymb,amsthm}
\usepackage{booktabs}
\usepackage[hidelinks]{hyperref}
\usepackage{enumitem}
\usepackage{url}
\usepackage{cleveref}

% Theorem environments
\theoremstyle{plain}
\newtheorem{theorem}{Theorem}[section]
\newtheorem{lemma}[theorem]{Lemma}
\newtheorem{corollary}[theorem]{Corollary}
\newtheorem{proposition}[theorem]{Proposition}
\theoremstyle{definition}
\newtheorem{definition}[theorem]{Definition}
\newtheorem{axiom}[theorem]{Axiom}
\newtheorem{assumption}[theorem]{Assumption}
\newtheorem{example}[theorem]{Example}
\theoremstyle{remark}
\newtheorem{remark}[theorem]{Remark}

\newcommand{\defeq}{\vcentcolon=}
\newcommand{\Fed}{\mathrm{Fed}}
\newcommand{\Hz}{H^0}
\newcommand{\Ho}{H^1}
\newcommand{\Res}{\mathrm{Res}}

\title{The Categorical Structure of Federated Convergence:\\
Limits, Sheaves, and Minimal Coordination}

\author{
Dayna Blackwell\\
\small\href{mailto:dayna@blackwell-systems.com}{dayna@blackwell-systems.com}\\
\small Licensed under CC-BY-4.0.}

\begin{document}
\maketitle

\begin{abstract}
\setlength{\parskip}{0.6\baselineskip plus 2pt minus 1pt}
Normalization confluence gives coordination-free convergence within a single registry, and its
federated extension gives convergence across a network of registries connected by morphisms, under
two decidable regimes: acyclic networks, and monotone repair on cycles~\cite{blackwell2026nc,
blackwell2026federated}. Those regimes are stated as separate sufficient conditions. This paper
gives the structure underneath them and, with it, an answer to a sharper question: when must a
federation coordinate at all, and how little coordination suffices?

We show that a federation's convergent states are a limit and its federated normalizer is the
retraction onto that limit, so compositionality (a verified sub-federation collapses to an effective
registry) is a corollary of one structural fact rather than a per-topology argument. We then show
that convergence certificates form a sheaf whose gluing condition is exactly the source-determinacy
and validity-preservation hypotheses the federation already certifies, so the obstruction to a global
convergent state is cohomological: $\Hz$ is the convergent states and $\Ho$ is the holonomy of the
morphism cycles. The two convergence regimes are then one picture: an acyclic network has no cycles
and $\Ho$ vanishes, while a monotone cycle collapses the obstruction to a least fixed point. The
obstruction is computable as a loop composite over the shared subspace, implemented in the reference
engine as a cycle diagnostic that names an obstructing cycle and exhibits its orbit. Finally we show
that when the obstruction is nonzero, coordinating a cycle basis of the morphism network suffices to
converge, and in the information-preserving (invertible) fragment this coordinated core is minimal:
a first-Betti-number many loops, localizing the global coordination verdict of CALM and invariant
confluence into a minimal, mixed-consistency partition. The structural core (the limit, the
retraction, and compositionality) is mechanized axiom-free in Coq/Rocq.
\end{abstract}

\medskip\noindent\textbf{Keywords:} Federated Convergence, Normalization Confluence, Sheaf
Cohomology, Coordination Avoidance, CALM, Category Theory, Limits, Holonomy, Minimal Coordination,
Mechanized Proof

\clearpage

%========================================
\section{Introduction}
\label{sec:intro}

Normalization confluence certifies that a single registry converges: under well-founded compensation
and compensation commutativity, every processor that consumes the same events reaches the same valid
state regardless of order~\cite{blackwell2026nc}. Its federated extension certifies that a network of
such registries, connected by directed morphisms that carry one registry's authority over another's
shared state, converges as a whole, in two regimes: any acyclic network whose morphisms preserve
validity, and any network (cyclic included) whose repair is monotone~\cite{blackwell2026federated}.
Both regimes are established by finite enumeration and are what the reference engine checks at build
time.

Those two regimes are presented as independent sufficient conditions, proved by separate arguments.
That leaves a structural question unanswered and a practical one unasked. The structural question:
is there one reason both regimes work, rather than two? The practical question: when a federation
falls into neither regime, the engine rejects it, but rejection is a blunt verdict. A designer wants
to know not merely that the network cannot converge coordination-free, but \emph{where} the conflict
is and \emph{how little} coordination would repair it. This paper answers both, and the answers turn
out to be the same object viewed twice.

\paragraph{The local-to-global view.}
A federation poses a local-to-global problem. Each registry is convergent on its own; each morphism,
in isolation, is a clean rule for propagating authority. The question is whether, once every morphism
constraint holds at once, there is a single global assignment the whole network agrees on, and
whether it is reached regardless of event order. This is the shape of a limit, and of a sheaf gluing
problem, and the obstruction to it is the shape of a cohomology class. Making that precise is what
unifies the regimes: an acyclic network has no cycles to obstruct gluing, and a monotone cycle has an
order that forces gluing even when a cycle is present. The same $\Ho$ that measures the obstruction
also, when it is nonzero, pinpoints the cycles responsible, which is exactly the diagnostic and the
minimal-coordination answers the practical question wants.

\paragraph{Scope, stated up front.}
The cleanest form of the cohomological classification holds in the \emph{invertible} fragment, where
each morphism transports its shared value by a bijection (a lossless relabeling: a copy, a
permutation, a bijective unit or schema map). There the transitions are group elements, $\Ho$ is
genuine \v{C}ech cohomology, and the minimal coordinated core is a cycle basis of exactly the first
Betti number. When morphisms collapse information (a clamp, a reset, a most-restrictive-wins merge,
the shape most compensation takes), the transitions are monoid elements rather than group elements,
classical $\Ho$ does not apply, and the correct statement is the operational one: a global convergent
state exists iff the loop composite of each cycle has a reachable fixed point. We therefore lead with
the general, operational results (the limit, the retraction, compositionality, and the loop-composite
diagnostic), which hold for arbitrary morphisms and are what the engine implements, and present the
cohomology as the sharp closed form on the invertible fragment. We state which real federations
inhabit that fragment (Section~\ref{sec:minimal}) rather than leaving it as ``the case where the math
is clean.''

\paragraph{Contributions.}
\begin{enumerate}[leftmargin=*]
\item \textbf{Federation as a limit, and compositionality for free
  (Sections~\ref{sec:limit}--\ref{sec:compose}).} A registry's valid set is the equalizer of its
  idempotent normalizer, a federation's convergent set is the equalizer of its shared-projection and
  resolver maps, and the federated normalizer is the idempotent retraction onto that limit. Because
  the normal forms are a limit, a convergent sub-federation collapses to an effective registry and
  the combined normalizer factors through it; compositionality is a corollary of one structural fact,
  not a per-topology theorem.
\item \textbf{Convergence certificates form a sheaf, and $R1/R2$ is its gluing axiom
  (Section~\ref{sec:sheaf}).} Local convergence certificates glue to a global one exactly when the
  subsystems agree as normalizers on their overlap. The two regimes that force agreement,
  single-writer authority and monotone overlap, are the federation's source-determinacy ($R1$) and
  validity-preservation ($R2$) hypotheses. The existing federation design is validated as the sheaf
  gluing check, not extended by a new mechanism.
\item \textbf{The obstruction is cohomological, and it unifies the two regimes
  (Section~\ref{sec:completion}).} $\Hz$ is the convergent states; $\Ho$ (in the invertible fragment)
  is the holonomy of the morphism cycles, generated by a cycle basis. $\Ho$ vanishes in exactly two
  ways, and they are the two convergence regimes: an acyclic network (no cycles) and a monotone cycle
  (a least fixed point). The dichotomy the base work states as two conditions is one invariant.
\item \textbf{A computable obstruction diagnostic (Section~\ref{sec:diagnostic}).} For a general
  (possibly non-invertible) cycle, a global convergent state exists iff the cycle's loop composite has
  a reachable fixed point, and the witness for its absence is the loop composite itself, computable
  over the shared subspace rather than the product. This is implemented in the reference engine as a
  cycle diagnostic that names the offending cycle and reports whether its repair settles or orbits.
\item \textbf{Minimal coordination (Section~\ref{sec:minimal}).} When the obstruction is nonzero,
  coordinating a cycle basis of the morphism network trivializes it, yielding a convergent
  implementation that is strongly consistent only on those loops and coordination-free everywhere
  else; in the invertible fragment this core is minimal, of size the first Betti number. This refines
  the global coordination verdict of CALM~\cite{hellerstein2010declarative,ameloot2013relational} and
  invariant confluence~\cite{bailis2014coordination} into a localized, minimal, mixed-consistency
  partition, and it upgrades the engine's blunt rejection of a nonconvergent cycle into an
  actionable ``coordinate these loops.''
\item \textbf{Mechanization.} The structural core (Contributions~1 and~2: the limit, the retraction
  for arbitrary acyclic federations, and compositionality) is mechanized axiom-free in Coq/Rocq,
  extending the machine-checked development of the base results.
\end{enumerate}

%========================================
\section{Related Work}
\label{sec:related}

\paragraph{Sheaf cohomology for distributed systems: the same tool, a different object.}
The use of sheaf cohomology to measure a local-to-global obstruction is not new, and we do not claim
it. Cellular-sheaf cohomology characterizes solvability of distributed decision \emph{tasks}, where
terminating solutions are the global sections and the cohomology encodes the obstruction to a
solution~\cite{riess2025sheaftasks}, in the combinatorial-topology lineage of distributed
computing~\cite{herlihy2013topology}; the same presheaf-and-obstruction template measures quantum
contextuality~\cite{abramsky2011sheaf}. Our object is different. Those works concern which one-shot
decision task a set of processes can solve under a failure or adversary model. We concern the
convergence of a \emph{federated normalization}: a network of registries, each with invariants and
compensation, whose morphisms must agree on shared state so that all processors reach the same valid
global normal form regardless of event order. Our presheaf is over the overlaps of a normalizer
structure, its sections are convergence certificates, $\Hz$ is the convergent states, and $\Ho$ is
the holonomy of the morphism cycles. To our knowledge no prior sheaf-cohomology work targets the
eventual-consistency or replicated-convergence obstruction, and none yields a coordination-minimality
result.

\paragraph{Coordination avoidance: a global verdict we localize.}
The CALM theorem states that a problem has a coordination-free, consistent implementation iff it is
monotone~\cite{hellerstein2010declarative, ameloot2013relational}, and invariant
confluence~\cite{bailis2014coordination} gives a related global avoidability criterion; recent work
sharpens the monotonicity boundary at the semantic level~\cite{completecalm2026,
coordfree2025}. These answer \emph{whether} coordination can be avoided, as a global property of the
whole problem. Our minimal-coordination result answers \emph{where} and \emph{how little}: it locates
the irreducible coordination in a cycle basis of the morphism network and shows the rest can remain
coordination-free, a localized mixed-consistency partition rather than a single yes-or-no. In the
information-preserving fragment the core is provably minimal. We are not aware of prior work that
localizes the coordination requirement to a cycle basis or bounds it by a Betti number.

\paragraph{Convergent replicated data and its algebra.}
Conflict-free replicated data types achieve convergence by restricting to operations that never need
repair; a state-based CRDT is the join-semilattice special case~\cite{shapiro2011crdt}, and the
monotone regime here coincides with it (this correspondence is machine-checked in the base
development). The well-founded-compensation regime is strictly broader, since it admits operations
that violate invariants and are repaired. Category-theoretic accounts of CRDT convergence exist at
the level of the merge algebra; they do not give a cohomological obstruction or a
coordination-minimality theorem. The lattice-theoretic and chaotic-iteration machinery behind the
monotone regime is classical~\cite{tarski1955, cousot1977}, as is the confluence core~\cite{newman1942}.

%========================================
\section{Background}
\label{sec:background}

We restate only what later sections use: the single-registry model (well-founded compensation and
compensation commutativity, giving unique normal forms via Newman's Lemma) and the federated model
(registry morphisms, the authority argument, and the single-source $M1$ and multi-source $R1/R2$
conditions). Full statements and proofs are in the base papers~\cite{blackwell2026nc,
blackwell2026federated}. [To be written: a self-contained half-page recap.]

%========================================
\section{Normal Forms Are a Limit}
\label{sec:limit}

[To be written from CATEGORICAL-STRUCTURE.md \S1--\S3. Lemma~0 (a registry's valid set is the
equalizer of its idempotent normalizer), Proposition~1 (the federated convergent set is the equalizer
of the shared-projection and resolver maps, a finite limit), and Theorem~1 (the federated normalizer
is the idempotent retraction onto that limit, order-independent). The retraction, the general acyclic
operator, and the limit characterization are mechanized in \texttt{coq/Categorical.v}.]

%========================================
\section{Compositionality for Free}
\label{sec:compose}

[To be written from CATEGORICAL-STRUCTURE.md \S5. Theorem~2 (a verified sub-federation collapses to
an effective registry; the flat and staged normalizers agree), and the assume-guarantee form with
input and output ports. The fold-append form of Theorem~2 is mechanized in \texttt{coq/Categorical.v}.]

%========================================
\section{Convergence Certificates Form a Sheaf}
\label{sec:sheaf}

[To be written from CATEGORICAL-STRUCTURE.md \S10. The presheaf of certificates, the non-naive gluing
counterexample, the precise gluing condition (agree as normalizers on the overlap), and the
identification of $R1 + R2$ with the gluing axiom.]

%========================================
\section{The Cohomological Completion}
\label{sec:completion}

[To be written from CATEGORICAL-STRUCTURE.md \S10.1. $\Hz$ as the convergent states; $\Ho$ as the
holonomy in the invertible fragment; the completion theorem; the cycle-basis generators; the two
routes to $\Ho = 0$ recovering the acyclic and monotone regimes; and the reduction of the general
non-invertible case to the loop-composite fixed-point condition.]

%========================================
\section{The Obstruction Diagnostic}
\label{sec:diagnostic}

[To be written from CATEGORICAL-STRUCTURE.md \S10. The loop composite as the computable witness,
bounded by the shared subspace; the settle-versus-orbit dichotomy; and the reference implementation
(\texttt{Federation.DiagnoseCycle}), with worked settle and orbit examples.]

%========================================
\section{Minimal Coordination}
\label{sec:minimal}

[To be written from CATEGORICAL-STRUCTURE.md \S10.2, incorporating the inhabitation analysis. The
minimal-coordination proposition; the coordination-free-majority plus coordinated-core decomposition;
the characterization of the invertible (lossless-relabeling) fragment and which real federations
inhabit it; the accept-with-coordination capability; and the positioning against CALM and invariant
confluence.]

%========================================
\section{Conclusion}
\label{sec:conclusion}

[To be written.]

%========================================
\begin{thebibliography}{99}

\bibitem{blackwell2026nc}
Dayna~Blackwell,
``Normalization Confluence for Registry-Governed Stream Processing,''
Zenodo, 2026.
\textsc{doi}: \href{https://doi.org/10.5281/zenodo.18671870}{10.5281/zenodo.18671870}.

\bibitem{blackwell2026federated}
Dayna~Blackwell,
``Normalization Confluence in Federated Registry Networks,''
Zenodo, 2026.
\textsc{doi}: \href{https://doi.org/10.5281/zenodo.18677400}{10.5281/zenodo.18677400}.

\bibitem{newman1942}
M.~H.~A. Newman,
``On theories with a combinatorial definition of `equivalence',''
\emph{Annals of Mathematics}, vol.~43, no.~2, pp.~223--243, 1942.

\bibitem{shapiro2011crdt}
M.~Shapiro, N.~Pregui\c{c}a, C.~Baquero, and M.~Zawirski,
``A comprehensive study of convergent and commutative replicated data types,''
\emph{INRIA Research Report RR-7506}, 2011.

\bibitem{hellerstein2010declarative}
J.~M.~Hellerstein,
``The declarative imperative: experiences and conjectures in distributed logic,''
\emph{SIGMOD Record}, vol.~39, no.~1, pp.~5--19, 2010.

\bibitem{ameloot2013relational}
T.~J.~Ameloot, F.~Neven, and J.~Van~den~Bussche,
``Relational transducers for declarative networking,''
\emph{Journal of the ACM}, vol.~60, no.~2, pp.~1--38, 2013.

\bibitem{bailis2014coordination}
P.~Bailis, A.~Fekete, M.~J.~Franklin, A.~Ghodsi, J.~M.~Hellerstein, and I.~Stoica,
``Coordination avoidance in database systems,''
\emph{Proc.\ VLDB Endowment}, vol.~8, no.~3, pp.~185--196, 2014.

\bibitem{abramsky2011sheaf}
S.~Abramsky and A.~Brandenburger,
``The sheaf-theoretic structure of non-locality and contextuality,''
\emph{New Journal of Physics}, vol.~13, 113036, 2011.

\bibitem{riess2025sheaftasks}
Authors of arXiv:2503.02556,
``A Sheaf-Theoretic Characterization of Tasks in Distributed Systems,''
arXiv:2503.02556, 2025.

\bibitem{herlihy2013topology}
M.~Herlihy, D.~Kozlov, and S.~Rajsbaum,
\emph{Distributed Computing Through Combinatorial Topology},
Morgan Kaufmann, 2013.

\bibitem{completecalm2026}
Authors of arXiv:2602.09435,
``Complete CALM: A Coordination Criterion for Specifications,''
arXiv:2602.09435, 2026.

\bibitem{coordfree2025}
Authors of arXiv:2504.01141,
``A Preliminary Model of Coordination-free Consistency,''
arXiv:2504.01141, 2025.

\bibitem{tarski1955}
A.~Tarski,
``A lattice-theoretical fixpoint theorem and its applications,''
\emph{Pacific Journal of Mathematics}, vol.~5, no.~2, pp.~285--309, 1955.

\bibitem{cousot1977}
P.~Cousot and R.~Cousot,
``Abstract interpretation: a unified lattice model for static analysis of programs by construction or approximation of fixpoints,''
\emph{Proc.\ 4th ACM SIGACT-SIGPLAN Symp.\ on Principles of Programming Languages (POPL)}, pp.~238--252, 1977.

\end{thebibliography}

\end{document}
